\begin{abstract}
Recovering the luminosity distance of a strongly lensed gravitational-wave (GW) source requires the absolute magnification of its images, which GW data alone cannot measure. We quantify how accurately and how reliably it can be inferred by combining the GW signals with an electromagnetic model of the lens galaxy. We simulate $8\times10^5$ lensed binary black holes from the GWTC-4.0 population, four galaxy-lens models from the singular isothermal sphere (SIS) to the elliptical power law with external shear, and O5-like and Einstein Telescope (ET) detector networks. We assume an identified host, reconstructed image positions of the GW source and no external convergence. Image geometry, the GW magnification ratio and the image arrival order---observables that contain no distance scale---give median errors of 3, 7, 10 and 19\% for the four lens models with ET and 10--27\% with O5, conditional on the adopted lens populations. Time delays and redshifts reduce the errors further only when a cosmology is assumed. Errors grow near caustics, where the 90\% intervals contain the true value for only about half of the images with $\mu>20$. An SIS prior gives biased estimates with underestimated uncertainties for realistic lenses, whereas calibrating on elliptical power-law lenses gives coverage of 0.89--0.99 for all four lens models. Parameter estimation of 120 lensed pairs from simulated detector data gives smaller uncertainties than our simplified GW measurement model. Conventional Bayesian lens reconstructions of the same systems reach 8\% (ET) and 14\% (O5) when the lens family is known, comparable to the trained estimators, but their nominal 90\% intervals contain only about 80\% of the ET truths. The radial mass profile of the lens and the positions of the GW images in the host are the main limitations.
\end{abstract}
